\subsection{Regularity at every collision place}
The preceding exception is excluded by a stronger
regularity comparison. The following operator argument
requires neither commutation with the actual action
nor restrictions on the base spectrum.

\begin{finding}[\fid{F55}, Part 2: the all-places inserted-trace gate]
\label{f:allplacestrace}
Let $\mathcal L/\C(v)$ be a finite coefficient field
for an algebraic matrix family $F(v)$ on the actual
rank-$64$ tensor module.
Assume regular evaluations at $b>1$ and all
sufficiently large $b^n$, with $F_b$ invertible.
If
\[
 \Tr(X_{b^n}F_{b^n})=\Tr(X_bF_b^n)
\]
for all sufficiently large $n$, then $F$ cannot be
regular on the actual pulled-back tensor lattice
at every place of $\mathcal L$ above $v=1$.
Here $X$ is the actual marked central operator,
with $X(1)=0$ on that lattice.

No commutation of $F_b$ with $X_b$, or of $F$ with
the actual right action, is needed.
The result allows arbitrary Jordan form and imposes
no common-radius, weight, or positivity hypothesis.
\end{finding}

\begin{proof}
We prove a general inserted-trace principle.

\par\smallskip\noindent\textbf{Polynomial identities propagate to index zero.}\quad
Let $W$ be a nonzero finite-dimensional complex vector
space, let $F_0$ be invertible and $X_0$ arbitrary, and set
\[
 S_n=\Tr(X_0F_0^n)\qquad(n\ge0).
\]
For any polynomial
$Q(v,Y)=\sum_{i,j}q_{ij}v^iY^j$, if $Q(b^n,S_n)$
vanishes for all sufficiently large $n$, then
\begin{equation}\label{eq:insertedtracezeroindex}
 Q(1,\Tr X_0)=0.
\end{equation}
For the nonzero polynomial case, on the finite direct sum
\[
 W_Q=\bigoplus_{q_{ij}\ne0}W^{\otimes j}
\]
put
\[
 T_Q=\bigoplus_{q_{ij}\ne0}b^iF_0^{\otimes j},
 \qquad
 U_Q=\bigoplus_{q_{ij}\ne0}q_{ij}X_0^{\otimes j},
\]
with zeroth tensor power equal to $\C$ and its
operators equal to $1$.
The operator $T_Q$ is invertible.
Tensor-product trace gives
\[
 z_n:=\Tr(U_QT_Q^n)=Q(b^n,S_n).
\]
No commutation is used: each block product is
$q_{ij}b^{in}(X_0F_0^n)^{\otimes j}$.

Write the characteristic polynomial of $T_Q$ as
$T^r+c_{r-1}T^{r-1}+\cdots+c_0$, with $c_0\ne0$.
Cayley--Hamilton gives
\[
 z_{n+r}+c_{r-1}z_{n+r-1}+\cdots+c_0z_n=0
 \qquad(n\ge0).
\]
Eventual vanishing therefore propagates backwards
one index at a time to $z_0=0$, proving
\eqref{eq:insertedtracezeroindex}.
This includes Jordan terms without choosing a
Jordan basis or asserting effectiveness of their traces.

\par\smallskip\noindent\textbf{A polynomial with prescribed reduction.}\quad
Let $f\in\mathcal L$ be regular with common residue
$c\in\C$ at every place above $v=1$.
Put $O=\C[v]_{(v-1)}$ and let $B$ be its integral closure
in $\mathcal L$.
It is finite over $O$ by the characteristic-zero
finite-normalization theorem
(Stacks Lemma~10.161.8,
\href{https://stacks.math.columbia.edu/tag/032L}{tag~032L}),
and finite free since it is torsion-free over this DVR
(Stacks Lemma~15.22.11(3),
\href{https://stacks.math.columbia.edu/tag/0AUW}{tag~0AUW}).
Its maximal localizations are the DVR places above one,
by lying over and the dimension-one normal-local criterion
(Stacks Lemmas~10.36.17 and~10.119.7,
\href{https://stacks.math.columbia.edu/tag/00GQ}{tag~00GQ}
and
\href{https://stacks.math.columbia.edu/tag/00PD}{tag~00PD}).

Regularity in all these localizations puts $f$ in $B$.
Indeed, if the denominator ideal
\[
 I=\{a\in B:af\in B\}
\]
were proper, choose a maximal ideal $\mathfrak m$
containing it. Membership $f\in B_{\mathfrak m}$
would give $a\notin\mathfrak m$ with $af\in B$,
a contradiction.

The finite $\C$-algebra $\overline B=B/(v-1)B$
has residue $\overline f=c$ at every maximal ideal.
Multiplication by $\overline f-c$ cannot have a
nonzero eigenvalue $\lambda$.
If it did, multiplication by $\overline f-c-\lambda$
would be singular, so this element would be a nonunit
and lie in a maximal ideal.
Its residue there is $-\lambda\ne0$, a contradiction.
Here multiplication is bijective exactly when its
multiplying element is a unit, and all residue fields
of $\overline B$ are $\C$.
Thus multiplication by $\overline f-c$ is nilpotent.

Choose an $O$-basis of $B$ and let $H(v,Y)\in O[Y]$
be the monic characteristic polynomial of multiplication
by $f$. Its degree $d=[\mathcal L:\C(v)]$ is positive,
and
\[
 H(1,Y)=(Y-c)^d.
\]
Cayley--Hamilton applied to $1\in B$ gives $H(v,f)=0$.
Clear its rational coefficient denominators by one
polynomial $a(v)$ with $a(1)\ne0$.
Then $Q=aH\in\C[v,Y]$ satisfies
\begin{equation}\label{eq:allplacespolynomial}
 Q(v,f)=0,\qquad Q(1,Y)=a(1)(Y-c)^d.
\end{equation}

Consequently, whenever the selected evaluations of
this same algebraic element satisfy
\[
 f_{b^n}=\Tr(X_0F_0^n)
\]
for all sufficiently large $n$,
\eqref{eq:insertedtracezeroindex} and
\eqref{eq:allplacespolynomial} force
\begin{equation}\label{eq:allplaceszerothtrace}
 c=\Tr X_0.
\end{equation}

\par\smallskip\noindent\textbf{The actual marked element.}\quad
Apply this principle to
\[
 f(v)=\Tr_M(X(v)F(v)),\qquad X_0=X_b,\quad F_0=F_b.
\]
All-places regularity of $F$ and the actual identity
$X(1)=0$ make $f$ regular with common residue $0$.
But $X_b$ acts by $\Delta(b)$ on its rank-eight
marked sector and by zero on the others, so
\[
 \Tr X_b=8\Delta(b)\ne0\qquad(b>1).
\]
This contradicts \eqref{eq:allplaceszerothtrace}.
Neither invariance of that sector under $F_b$ nor
commutation with the right action was used.
\end{proof}

The same principle covers finite alternating inserted
traces: use a direct-sum base operator and put the
alternating signs into the insertion.
Zero residue at every place then requires
$\sum_j(-1)^j\Tr X_j=0$.
If that zeroth inserted Euler trace is zero,
this obstruction says nothing; it does not turn
a virtual trace into an effective spectrum.

Base invertibility is essential.
For $F_0=0$ and $X_0=I$ on a nonzero space,
the inserted traces vanish for all $n\ge1$ but their
zeroth trace is the dimension.
Actual finite Weil spaces have invertible Frobenius.
If $X$ specializes only to a nilpotent, one must
separately verify that $\Tr(XF)$ has zero residue;
commutation is one sufficient condition in that variant.

The explicit model in Part~1 has poles in its operator
at $t=+1$ and therefore fails precisely the all-places
regularity hypothesis.
Its regularity at $t=-1$, actual right-action
compatibility and two trace laws remain valid.
All-places regularity is extra comparison data;
it is not supplied automatically by an algebraic
coefficient scheme or a perverse realization.

% Separate paragraph for Part III.
\paragraph{\fid{F55}, Part 2, strengthened, in Grothendieck's language.}
The normalization $B$ is a finite free algebra over
the local parameter ring, and multiplication by $f$
is an actual endomorphism whose characteristic
polynomial records the special-fiber scalar reduction.
The underlying affine module language is EGA~I,
Theorem~1.3.7, p.~86, Corollary~1.3.8, p.~87,
and Theorem~1.4.1, pp.~90--92 \cite{EGA}.
The exact normalization, freeness and place statements
are the primary Stacks lemmas cited in the proof.

The polynomial trace identity is represented by
specified finite direct sums and tensor powers
of endomorphisms. The alternating extension uses
an inserted Euler trace, not an effective-spectrum
or $\gamma$-positivity inference.
The established representation/Grothendieck framework
anchors remain SGA~3, I, 4.7.2--3 and
SGA~6, V, \S\S2--3, with cote~90, scan~17 and
cote~133, scan~18 \cite{SGA3,SGA6,Archive}.
The proof concerns this added all-places regularity
comparison with the actual lattice; it does not
identify a canonical basis or supply the independent
pointwise coefficient estimates.

